\centerline{\bf \Huge Геометрия 4}
\centerline{\bf \Huge Вокруг ортоцентра I}

\textbf{Основные свойства.}

(1) Направления из вершины треугольника на ортоцентр и центр описанной окружности симметричны относительно биссектрисы соответствующего угла треугольника;

(2) Точки, симметричные ортоцентру относительно сторон треугольника, лежат на описанной окружности данного треугольника;

(3) Точки, симметричные ортоцентру относительно середин сторон треугольника, лежат на описанной окружности данного треугольника, причем являются диаметрально противоположными к вершинам треугольника;

(4) Высоты треугольника являются биссектрисами ортотреугольника, то есть треугольника, вершинами которого являются основания высот исходного треугольника.

(5) \textbf{Окружность Эйлера (девяти точек))}. $A B C$ --- треугольник.

$A_1, B_1, C_1$ --- основания высот треугольника $A B C$ из вершин $A, B, C$ соответственно. $H$ ортоцентр.

$A_2, B_2, C_2$ - середины сторон $B C, A C, B A$ соответственно.

$A_3, B_3, C_3$ - середины отрезков $H A, H B, H C$ соответственно.

$A_1, B_1, C_1, A_2, B_2, C_2, A_3, B_3, C_3$ лежат на одной окружности.

Сперва докажите все эти свойства и только потом приступайте к решению задач. 

\newpage

\hspace{3mm}

\problem{1}
Предположим, что две вершины треугольника и его описанная окружность зафиксированы, а третья вершина движется по этой описанной окружности. По какой траектории будет двигаться ортоцентр треугольника? (Не забудьте, что вершина может двигаться по двум разным дугам!)

\problem{2}
Обозначим за $O$ и $H$ центр описанной около треугольника $A B C$ окружности и его ортоцентр. Докажите, что

(a) $\angle B A H=\angle C A O$

(b) Расстояние от $O$ до $B C$ равно половине отрезка $A H$.

\problem{3}
В остроугольном треугольнике $A B C$ угол при вершине $A$ равен $60^{\circ}$. Известно, что $B C=1$. Найдите длину отрезка $B_1 C_1$, если $B_1$ и $C_1$ это основания высот треугольника, опущенных из вершин $B$ и $C$.

\problem{4}
Пусть $H$ - ортоцентр остроугольного треугольника $A B C$. Точка $A^{\prime}$ выбирается на описанной около этого треугольника окружности $\Omega$ так, что $A A^{\prime}$ - диаметр $\Omega$. Докажите, что четырехугольник $B H C A^{\prime}$ - параллелограмм.

\problem{5}
Треугольник $A B C$ с ортоцентром $H$ вписан в окружность $\omega$. Окружность, построенная на отрезке $A H$ как на диаметре, пересекает $\omega$ в точке $S \neq A$. Докажите, что прямая $S H$ делит отрезок $B C$ пополам.

\problem{6}
Докажите, что середины трех отрезков, соединяющих ортоцентр с вершинами треугольника, лежат на окружности Эйлера этого треугольника.

\problem{7}
Четырехугольник $A B C D$ вписанный. Докажите, что ортоцентры треугольников $A B C, B C D$, $C D A$ и $D A B$ образуют четырехугольник, равный исходному.

\problem{8}
\textbf{Точка Анти-Штейнера.} Прямая $\ell$ проходит через ортоцентр треугольника $A B C$. Докажите, что три прямые, симметричные прямой $\ell$ относительно сторон треугольника, пересекаются в одной точке.